\documentclass{article}
\usepackage{amsmath, amssymb, amsthm}
\usepackage{hyperref}
\usepackage{geometry}
\geometry{margin=1in}

\title{Erd\H{o}s Problem \#1055}
\date{Last updated: 2025-09-28}
\author{Source: erdosproblems.com}

\begin{document}
\maketitle

\noindent\textbf{Status:} OPEN \quad \textbf{No prize}

\noindent\textbf{Tags:} number theory, primes

\medskip

\noindent\textbf{Problem Statement:}

A prime $p$ is in class $1$ if the only prime divisors of $p+1$ are $2$ or $3$. In general, a prime $p$ is in class $r$ if every prime factor of $p+1$ is in some class $\leq r-1$, with equality for at least one prime factor. Are there infinitely many primes in each class? If $p_r$ is the least prime in class $r$, then how does $p_r^{1/r}$ behave?




A classification due to Erd\H{o}s and Selfridge. It is easy to prove that the number of primes $\leq n$ in class $r$ is at most $n^{o(1)}$. The sequence $p_r$ begins $2,13,37,73,1021$ ( A005113 in the OEIS). Erd\H{o}s thought $p_r^{1/r}\to \infty$, while Selfridge thought it quite likely to be bounded. A similar question can be asked replacing $p+1$ with $p-1$. This is problem A18 in Guy's collection \cite{Gu04}.




References

[Gu04] Guy, Richard K., Unsolved problems in number theory . (2004), xviii+437.

\noindent\textbf{Related OEIS sequences:} A005113


\bigskip
\noindent\small{Source: \url{https://www.erdosproblems.com/1055}}
\end{document}
